\documentclass{stp}

\begin{document}

\stptitle{
  вид        = лабораторная,
  номер      = 1,
  дисциплина = {ИТиВП},
  студент    = {А. С. Образцов},
  группа     = {320601},
  преподаватель = {Н. В. Хаджинова},
}

\section{Цель работы}

Целью данной лабораторной работы является приобретение практических
навыков разработки серверного приложения на платформе Node.js с
использованием фреймворка Express.js. В процессе выполнения работы
изучаются принципы построения серверной части веб-приложения,
организация маршрутизации HTTP-запросов и формирование ответов
клиенту. В рамках лабораторной работы разрабатывается REST API для
курсового проекта — сервиса планирования задач по методу GTD (Getting
Things Done) с контекстами и напоминаниями.

Основной сущностью разрабатываемого приложения является задача
(task). Каждая запись содержит заголовок, контекст выполнения
(например, @calls, @computer, @errands),
статус в рамках методологии GTD (входящая, следующее действие, в
ожидании, когда-нибудь, выполнена), срок выполнения, время
напоминания и произвольные заметки. Для работы с коллекцией задач
реализованы CRUD-операции: получение списка всех задач с возможностью
фильтрации по контексту и статусу, получение отдельной задачи по
идентификатору, создание новой задачи, полное обновление существующей
и удаление записи. Дополнительно для каждого маршрута реализована
обработка запроса OPTIONS, возвращающего список разрешённых методов.

Метод GET используется для получения данных, POST — для создания
новой задачи, PUT — для полного изменения существующей записи, DELETE
— для её удаления, OPTIONS — для получения перечня методов,
поддерживаемых ресурсом. Идентификатор задачи передаётся как параметр
маршрута через req.params, фильтры контекста и статуса — как
параметры строки запроса через req.query, а данные для
создания и изменения задачи передаются в теле запроса в формате JSON.
Для разбора JSON-данных используется middleware express.json().

При успешном выполнении запросов сервер возвращает статусы 200~OK при
получении и обновлении данных, 201~Created при создании новой задачи
и 204~No~Content при успешном удалении. При передаче некорректных
данных используется статус 400~Bad~Request, а при обращении к
несуществующей задаче — 404~Not~Found. Предусмотрен также глобальный
обработчик непредвиденных ошибок сервера со статусом 500~Internal
Server Error.

На данном этапе задачи хранятся во временном массиве в памяти
приложения без использования базы данных. Такой подход позволяет
сосредоточиться на изучении принципов работы Node.js, Express.js и
REST API. В дальнейшем сервис расширяется подключением базы данных,
системой напоминаний, аутентификацией пользователей и интеграцией с
клиентской частью курсового проекта.

\section{Краткие теоретические сведения}

Node.js представляет собой среду исполнения JavaScript вне браузера,
построенную на движке V8. Событийно-ориентированная модель и
неблокирующий ввод-вывод позволяют Node.js эффективно обрабатывать
большое количество одновременных подключений при минимальных
накладных расходах, что важно для сервиса, синхронизирующего задачи и
напоминания пользователя в реальном времени.

Express.js — минималистичный фреймворк для Node.js, предоставляющий
инструменты маршрутизации и промежуточного программного обеспечения
(middleware). Маршрут связывает HTTP-метод и URL-путь с
функцией-обработчиком, принимающей объект запроса (req) и
объект ответа (res). Механизм express.Router()
позволяет группировать связанные маршруты в отдельный модуль и
подключать его к приложению по общему префиксу.

Архитектурный стиль REST (Representational State Transfer) определяет
правила построения API поверх протокола HTTP. Основными принципами
являются: клиент-серверная архитектура, единообразие интерфейса (URL
идентифицирует ресурс, метод — действие над ним), отсутствие
состояния на сервере между запросами (stateless), кэшируемость
ответов и возможность построения многоуровневой архитектуры.

Промежуточное программное обеспечение (middleware) — функции,
выполняемые в цепочке до передачи запроса конечному обработчику
маршрута. Функция express.json() разбирает тело запроса из
формата JSON в объект JavaScript и помещает результат в
req.body. Обработчик ошибок также является middleware, но
принимает четыре аргумента: err, req, res
и next, что позволяет Express отличать его от обычных
маршрутов.

Метод HTTP OPTIONS предназначен для запроса сведений о том, какие
методы поддерживаются конкретным ресурсом, без выполнения какого-либо
действия над ним. Сервер отвечает на такой запрос статусом 204~No
Content и заголовком Allow, перечисляющим допустимые методы.
Этот механизм используется браузерами при выполнении так называемых
предварительных (preflight) запросов CORS для составных
кросс-доменных обращений.

Формат передачи данных JSON — текстовый формат обмена данными,
основанный на синтаксисе объектов JavaScript, являющийся стандартом
де-факто для REST API. Сервер принимает JSON в теле запроса и
возвращает JSON в теле ответа, устанавливая заголовок
Content-Type: application/json. Каждый ответ сопровождается
статусным кодом: 200 означает успешное выполнение, 201 — создание
ресурса, 204 — успешное выполнение без тела ответа, 400 — ошибку в
запросе клиента, 404 — отсутствие запрошенного ресурса, 500 —
внутреннюю ошибку сервера.

\section{Ход работы}

\subsection{Инициализация проекта и создание сервера}

В каталоге проекта выполнена команда npm init, после чего
установлен фреймворк Express.js и инструмент nodemon для
автоматического перезапуска сервера при изменении кода. В файле
package.json добавлены скрипты start и dev
для обычного запуска сервера и запуска в режиме разработки
соответственно.

Точкой входа приложения является файл server.js. В нём
подключается Express, создаётся экземпляр приложения, регистрируется
middleware express.json(), задаётся порт~3000 и запускается
HTTP-сервер. Логика приложения разделена по каталогам
routes, controllers и models: маршруты
описывают только соответствие метода и пути обработчику, контроллер
содержит бизнес-логику и валидацию, а модель инкапсулирует хранение
задач во временном массиве.

Ниже приведён фрагмент точки входа приложения:

\begin{stplisting}
const express = require('express');
const tasksRouter = require('./routes/tasks.routes');
const { notFoundHandler, errorHandler }
  = require('./middleware/errorHandler');

const app = express();
const PORT = process.env.PORT || 3000;

app.use(express.json());
app.use('/tasks', tasksRouter);

app.use(notFoundHandler);
app.use(errorHandler);

app.listen(PORT, () =>
  console.log('Server running on :' + PORT));
\end{stplisting}

После выполнения команды npm run dev nodemon
запускает сервер, и в консоли выводится сообщение о том, что
приложение слушает адрес http://localhost:3000.

\subsection{Реализация REST API для задач GTD}

В памяти сервера создан массив tasks, содержащий тестовые
задачи. Каждая задача содержит идентификатор id, заголовок
title, контекст context, статус status,
срок dueDate, время напоминания reminderAt и
заметки notes. Значение идентификатора для новых записей
выдаётся счётчиком nextId.

Для каждого маршрута коллекции и отдельного ресурса, помимо основных
методов, зарегистрирован обработчик запроса OPTIONS, возвращающий
заголовок Allow со списком допустимых методов:

\begin{stplisting}
function allowCollection(req, res) {
  res.set('Allow', 'GET, POST, OPTIONS');
  res.status(204).send();
}
function allowItem(req, res) {
  res.set('Allow', 'GET, PUT, DELETE, OPTIONS');
  res.status(204).send();
}

router.get('/', tasksController.list);
router.post('/', tasksController.create);
router.options('/', allowCollection);

router.get('/:id', tasksController.getOne);
router.put('/:id', tasksController.update);
router.delete('/:id', tasksController.remove);
router.options('/:id', allowItem);
\end{stplisting}

Проверка тела запроса выполняется функцией validateTaskInput,
которая контролирует наличие заголовка задачи и допустимость
контекста и статуса из фиксированного набора значений:

\begin{stplisting}
function validateTaskInput(body) {
  const errors = [];
  if (!body.title || !body.title.trim()) {
    errors.push('Поле "title" обязательно');
  }
  if (!ALLOWED_CONTEXTS.includes(body.context)) {
    errors.push('Недопустимый контекст задачи');
  }
  return errors;
}
\end{stplisting}

Список всех задач поддерживает необязательную фильтрацию по
контексту и статусу через параметры строки запроса:

\begin{stplisting}
function list(req, res) {
  let result = tasksModel.getAll();
  const { context, status } = req.query;
  if (context) {
    result = result.filter(
      (t) => t.context === context);
  }
  if (status) {
    result = result.filter(
      (t) => t.status === status);
  }
  res.status(200).json(result);
}
\end{stplisting}

Для обращений к неизвестным маршрутам зарегистрировано завершающее
middleware, возвращающее 404~Not~Found. Глобальный обработчик ошибок
принимает четыре аргумента и различает две ситуации: некорректный
JSON в теле запроса (статус 400) и непредвиденную ошибку сервера
(статус 500):

\begin{stplisting}
function errorHandler(err, req, res, next) {
  console.error(err.stack);
  if (err.type === 'entity.parse.failed') {
    return res.status(400).json({
      error: 'Некорректный JSON в запросе',
    });
  }
  res.status(err.status || 500).json({
    error: err.message || 'Ошибка сервера',
  });
}
\end{stplisting}

\subsection{Тестирование API}

Тестирование выполнено с помощью утилиты curl при запущенном сервере
на порту~3000. Для каждого маршрута проверены HTTP-метод, URL,
возвращаемый статус и тело JSON-ответа. В запросах POST и PUT тело
передаётся с заголовком Content-Type: application/json.

Получение одной задачи по идентификатору выполняется запросом
GET~/tasks/3:

\begin{stplisting}
curl -s http://localhost:3000/tasks/3

{
  "id": 3,
  "title": "Сверстать отчёт по лабораторной",
  "context": "@computer",
  "status": "waiting",
  "dueDate": "2026-09-20",
  "reminderAt": "2026-09-19T18:00:00",
  "notes": "Ждём скриншоты из Postman",
  "createdAt": "2026-09-15T10:10:00.000Z"
}
                                 HTTP 200 OK
\end{stplisting}

Создание новой задачи проверено запросом POST~/tasks с
телом, содержащим заголовок, контекст, статус и срок выполнения:

\begin{stplisting}
curl -X POST http://localhost:3000/tasks \
  -H "Content-Type: application/json" \
  -d '{"title":"Сдать курсовой проект",
       "context":"@computer","status":"next",
       "dueDate":"2026-12-01"}'

{
  "id": 4,
  "title": "Сдать курсовой проект",
  "context": "@computer",
  "status": "next",
  "dueDate": "2026-12-01",
  "reminderAt": null,
  "notes": "",
  "createdAt": "2026-09-16T19:16:57.390Z"
}
                            HTTP 201 Created
\end{stplisting}

Сервер автоматически присвоил задаче новый идентификатор и вернул
статус 201~Created вместе с созданным объектом.

Полное обновление задачи проверено запросом PUT~/tasks/2: в
теле запроса переданы новые значения заголовка, контекста и статуса:

\begin{stplisting}
curl -X PUT http://localhost:3000/tasks/2 \
  -H "Content-Type: application/json" \
  -d '{"title":"Купить продукты для дома",
       "context":"@errands","status":"next"}'

{
  "id": 2,
  "title": "Купить продукты для дома",
  "context": "@errands",
  "status": "next",
  "dueDate": null,
  "reminderAt": null,
  "notes": "",
  "createdAt": "2026-09-15T10:05:00.000Z"
}
                                 HTTP 200 OK
\end{stplisting}

Удаление задачи проверено запросом DELETE~/tasks/2. При
успешном выполнении сервер вернул статус 204~No~Content без тела
ответа, а повторный запрос к тому же идентификатору — статус
404~Not~Found:

\begin{stplisting}
curl -i -X DELETE http://localhost:3000/tasks/2
HTTP/1.1 204 No Content

curl -X DELETE http://localhost:3000/tasks/2
{
  "error": "Задача с id=2 не найдена"
}
                            HTTP 404 Not Found
\end{stplisting}

Обращение к несуществующей задаче GET~/tasks/999 также
возвращает статус 404~Not~Found с описанием ошибки в теле ответа:

\begin{stplisting}
curl -s http://localhost:3000/tasks/999
{
  "error": "Задача с id=999 не найдена"
}
                            HTTP 404 Not Found
\end{stplisting}

Передача некорректных данных в POST-запросе (пустой заголовок и
отсутствующий контекст) приводит к статусу 400~Bad~Request с
перечислением ошибок валидации:

\begin{stplisting}
curl -X POST http://localhost:3000/tasks \
  -H "Content-Type: application/json" \
  -d '{"title":""}'

{
  "error": "Некорректные данные запроса",
  "details": [
    "Поле \"title\" обязательно и должно",
    "быть непустой строкой",
    "Поле \"context\" обязательно, один из:",
    "@home, @work, @calls, @errands,",
    "@computer"
  ]
}
                             HTTP 400 Bad Request
\end{stplisting}

Наконец, для проверки обработки запроса OPTIONS выполнены запросы к
коллекции и к отдельному ресурсу — сервер возвращает статус
204~No~Content и заголовок Allow с перечнем поддерживаемых
методов:

\begin{stplisting}
curl -i -X OPTIONS http://localhost:3000/tasks
HTTP/1.1 204 No Content
Allow: GET, POST, OPTIONS

curl -i -X OPTIONS http://localhost:3000/tasks/1
HTTP/1.1 204 No Content
Allow: GET, PUT, DELETE, OPTIONS
\end{stplisting}

Проведённое тестирование показало, что все реализованные маршруты
корректно выполняют CRUD-операции над задачами и возвращают
соответствующие HTTP-статусы как при успешном выполнении запроса, так
и при ошибочных данных или обращении к несуществующему ресурсу.

\subsection{Таблица маршрутов}

Полный перечень реализованных маршрутов приведён в
таблице~\ref{tab:routes}.

\begin{stptable}
  \caption{Маршруты разработанного REST API}\label{tab:routes}
  \begin{tabularx}{\linewidth}{|c|X|X|c|c|}\hline
    Метод & Маршрут & Назначение & Успешный статус & Ошибки\\\hline
    GET & /tasks & Список задач (с фильтрами context, status) & 200 OK & --\\\hline
    GET & /tasks/:id & Получение задачи по id & 200 OK & 400, 404\\\hline
    POST & /tasks & Создание новой задачи & 201 Created & 400\\\hline
    PUT & /tasks/:id & Полное обновление задачи & 200 OK & 400, 404\\\hline
    DELETE & /tasks/:id & Удаление задачи & 204 No Content & 400, 404\\\hline
    OPTIONS & /tasks & Список разрешённых методов коллекции & 204 No Content & --\\\hline
    OPTIONS & /tasks/:id & Список разрешённых методов ресурса & 204 No Content & --\\\hline
  \end{tabularx}
\end{stptable}

Представленная таблица отражает основные маршруты разработанного
REST API, их назначение, успешные HTTP-статусы и возможные коды
ошибок. Для операций создания и полного обновления задачи клиент
должен передавать данные в теле запроса в формате JSON; обязательными
являются заголовок задачи и контекст из фиксированного набора
значений, поле статуса при отсутствии принимается равным
inbox.

\section{Ответы на контрольные вопросы}

\begin{stpnum}
  \item Реализованы методы GET, POST, PUT, DELETE и OPTIONS. GET
    используется для получения списка задач и отдельной задачи по
    идентификатору, POST — для создания новой задачи, PUT — для
    полного обновления существующей записи, DELETE — для её
    удаления, OPTIONS — для получения списка методов, допустимых для
    ресурса.
  \item Параметры маршрута доступны через req.params: в
    данной работе идентификатор задачи читается как
    req.params.id в маршруте /tasks/:id. Параметры
    строки запроса доступны через req.query: в списке задач
    они используются для необязательной фильтрации по контексту
    (?context=@computer) и статусу (?status=next).
  \item Middleware express.json() разбирает тело входящего
    HTTP-запроса в формате JSON и помещает результат в объект
    req.body. Без этого middleware при отправке JSON в
    POST- и PUT-запросах свойство req.body осталось бы
    пустым или неопределённым.
  \item При успешном создании ресурса сервер возвращает статус
    201~Created вместе с созданным объектом. При успешном удалении
    возвращается статус 204~No~Content, поэтому тело ответа
    отсутствует.
  \item REST API — интерфейс взаимодействия клиента и сервера,
    построенный вокруг ресурсов и стандартных методов HTTP. К
    основным принципам REST относятся клиент-серверная архитектура,
    отсутствие состояния между запросами (stateless), единообразный
    интерфейс, кэшируемость ответов и возможность построения
    многоуровневой архитектуры.
  \item Глобальная обработка ошибок организуется через middleware с
    четырьмя параметрами (err, req, res,
    next), размещаемое после всех маршрутов приложения. В
    данной работе такой обработчик различает некорректный JSON в
    теле запроса (статус 400) и непредвиденную ошибку сервера
    (статус 500), предварительно записывая её в консоль.
  \item Для ошибок выбрана единая структура ответа
    \{ "error": "текст ошибки" \}, а для ошибок валидации
    дополнительно добавляется поле details с перечнем
    конкретных нарушений, например отсутствия обязательного поля или
    недопустимого значения контекста.
  \item Нужно дважды отправить DELETE-запрос с одним и тем же
    идентификатором. Первый запрос удаляет существующую задачу и
    возвращает статус 204~No~Content. Повторный запрос должен
    вернуть статус 404~Not~Found, поскольку такой задачи в массиве
    уже нет; это поведение проверено при тестировании.
  \item Инструмент nodemon отслеживает изменения файлов проекта
    и автоматически перезапускает Node.js-сервер после сохранения
    кода. Это ускоряет разработку, поскольку сервер не требуется
    останавливать и запускать вручную после каждого изменения.
  \item PUT предназначен для полного обновления ресурса: клиент
    передаёт все поля нового состояния задачи. PATCH используется
    для частичного изменения отдельных полей без передачи остальных.
    В работе реализован PUT, поскольку он указан в задании и
    соответствует полной замене записи о задаче.
  \item Для поддержки CORS необходимо установить пакет cors
    командой npm install cors, подключить его через
    const cors = require('cors') и зарегистрировать
    middleware app.use(cors()). Это потребуется при
    подключении клиентской части курсового проекта, работающей с
    другого origin.
  \item Массив используется для упрощения первой лабораторной работы
    и позволяет сосредоточиться на маршрутизации HTTP-запросов и
    CRUD-операциях. Недостаток такого подхода — потеря всех данных
    при перезапуске сервера. Переход на базу данных необходим, когда
    требуется постоянное хранение задач и напоминаний, связи между
    сущностями и одновременная работа нескольких пользователей.
  \item Для отправки JSON-данных необходимо установить заголовок
    Content-Type: application/json. Именно на этот заголовок
    реагирует middleware express.json() при разборе тела
    запроса.
  \item Пример команды: curl -X POST
    http://localhost:3000/tasks -H "Content-Type: application/json"
    -d '\{"title":"Позвонить в деканат","context":"@calls"\}'.
  \item Необходимо следить за выводом nodemon в терминале,
    проверять отсутствие ошибок после автоматического перезапуска и
    повторно выполнять контрольные запросы через curl, контролируя
    HTTP-статус и тело ответа для каждого изменённого маршрута.
  \item Необходимо установить Node.js, получить исходный код проекта
    из репозитория (ветка lab21), перейти в каталог проекта
    и выполнить npm install для установки зависимостей.
    После этого сервер запускается командой npm start или
    npm run dev; требуется также убедиться, что порт~3000
    свободен.
\end{stpnum}

\stpunnumbered{Заключение}

В ходе выполнения лабораторной работы разработано серверное
приложение на платформе Node.js с использованием фреймворка
Express.js. Предметной областью приложения является сервис
планирования задач по методу GTD с поддержкой контекстов и
напоминаний — часть курсового проекта. Изучены принципы построения
серверной части веб-приложения, маршрутизации HTTP-запросов и
формирования ответов клиенту.

Для взаимодействия с данными реализован REST API для коллекции
tasks: получение списка задач с фильтрацией по контексту и
статусу, получение отдельной задачи по идентификатору, создание,
полное обновление и удаление существующей записи. Для каждого
маршрута дополнительно реализована обработка запроса OPTIONS,
возвращающего заголовок Allow со списком поддерживаемых
методов. Обмен данными между клиентом и сервером осуществляется в
формате JSON.

На данном этапе задачи временно хранятся в массиве в памяти сервера,
что позволило сосредоточиться на изучении CRUD-операций, параметров
маршрута, строки запроса и тела запроса без подключения базы данных.
Реализована обработка некорректных запросов: при отсутствии
обязательных полей или недопустимом значении контекста сервер
возвращает статус 400~Bad~Request, при обращении к несуществующей
задаче — 404~Not~Found, при успешном создании — 201~Created, при
успешном удалении — 204~No~Content. Для непредвиденных ошибок
предусмотрен глобальный обработчик со статусом 500~Internal Server
Error.

Все реализованные маршруты протестированы с помощью утилиты curl.
Тестирование подтвердило корректную работу операций получения,
создания, обновления и удаления задач, а также запроса OPTIONS и
обработку ошибочных запросов. Разработанное приложение используется
как основа для последующих лабораторных работ и дальнейшего развития
курсового проекта, включая подключение базы данных, систему
напоминаний, аутентификацию пользователей и интеграцию с клиентской
частью.

\end{document}
